\begin{definitionbox}{Definition --- Open Ball}
\begin{definition}[Open Ball]\label{def:metric-open-ball}
Let \((X,d)\) be a metric space, let \(x\in X\), and let \(r>0\).  The
\emph{open ball} of radius \(r\) centered at \(x\) is
\[
  B_d(x;r)
  :=
  \{\,y\in X : d(y,x)<r\,\}.
\]
\end{definition}
\end{definitionbox}

\begin{remark*}[Standard quantified statement]
\[
\forall y\in X{:}\quad
y\in B_d(x;r)
\Longleftrightarrow
d(y,x)<r.
\]
\end{remark*}

\begin{remark*}[Interpretation]
An open ball collects exactly the points whose distance from the center is
strictly less than the chosen radius.  The strict inequality leaves the
boundary points out.
\end{remark*}

\begin{remark*}[Historical note]
This notation follows Definition~2.15 in Kainth's
\emph{A Comprehensive Textbook on Metric Spaces}
\cite{KainthComprehensiveMetricSpaces2023}.
\end{remark*}

\begin{dependencies}
\begin{itemize}
  \item \hyperref[def:metric-space]{Definition: Metric Space}
\end{itemize}
\end{dependencies}

\begin{definitionbox}{Definition --- Closed Ball}
\begin{definition}[Closed Ball]\label{def:metric-closed-ball}
Let \((X,d)\) be a metric space, let \(x\in X\), and let \(r>0\).  The
\emph{closed ball} of radius \(r\) centered at \(x\) is
\[
  \overline{B}_d(x;r)
  :=
  \{\,y\in X : d(y,x)\leq r\,\}.
\]
\end{definition}
\end{definitionbox}

\begin{remark*}[Standard quantified statement]
\[
\forall y\in X{:}\quad
y\in \overline{B}_d(x;r)
\Longleftrightarrow
d(y,x)\leq r.
\]
\end{remark*}

\begin{remark*}[Interpretation]
A closed ball uses the same center and radius as the corresponding open ball,
but it includes the boundary points at distance exactly \(r\).
\end{remark*}

\begin{dependencies}
\begin{itemize}
  \item \hyperref[def:metric-space]{Definition: Metric Space}
  \item \hyperref[def:metric-open-ball]{Definition: Open Ball}
\end{itemize}
\end{dependencies}
